%% parts/p03/qe-categoricity-games-supplement.tex
%% Separately attributed derivative of historical additions.
\noindent\textit{Reader note.} The following separately authored material
retains the nonoverlapping additions from the preserved historical QE draft in
their original order.  Its displayed citations and proof boundaries control
its claims; it supplements the preceding
\emph{Quantifier Elimination and Categoricity} chapter.

%% BEGIN AUDITED QE BODY SLICE 104-212
The characteristic is part of the theory: the union customarily called
``ACF'' without a fixed characteristic is not complete.  Real closed fields,
by contrast, have literal elimination in the \emph{ordered}-ring language;
Presburger arithmetic needs congruence or divisibility predicates before an
elimination theorem is literal.  These language boundaries are made explicit
in the two catalogs below \cite{Marker2002,Presburger1929}.

\begin{hbdefinition}{def:p03-definable-skolem-functions}{Definable Skolem functions}
A theory \(T\) has \emph{definable Skolem functions} if, for every formula
\(\varphi(\bar x,y)\), some formula \(\theta(\bar x,y)\) satisfies
\[
 T\entails\forall\bar x\Bigl(
   \exists y\,\varphi(\bar x,y)\to
   \exists!y\,[\varphi(\bar x,y)\land\theta(\bar x,y)]\Bigr).
\]
This is a property of a theory: \(\theta\) uniformly chooses a definable
witness on every nonempty fiber.  It is separated here because neither
quantifier elimination nor model completeness alone supplies such a choice;
the next catalog uses the distinction.
\index{definable Skolem functions}
\end{hbdefinition}

For the last complexity cell of Table~\ref{tab:p03-qe-catalog},
\(\mathrm{STA}(s,t,a)\) denotes the class accepted by alternating Turing
machines using space \(O(s(n))\), time \(O(t(n))\), and at most \(O(a(n))\)
alternations; an asterisk leaves that resource unrestricted.  This
metalinguistic abbreviation is needed to report Kozen's exact Boolean-algebra
classification and is not a synonym for \(\mathsf{PSPACE}\).

In both tables below, the \emph{Boolean-algebra signature} means a five-symbol
nonlogical signature with two nullary constants, two fresh binary operation
symbols, and one fresh unary complement symbol, all disjoint from the logical
alphabet and interpreted as \(0,1\), meet, join, and complement.

\begin{landscape}
\begingroup\tiny\setlength{\tabcolsep}{3pt}
\begin{longtable}{@{}p{1.15in}p{1.05in}p{.62in}p{.65in}p{.95in}p{2.64in}@{}}
\caption{First-order theories master table.  ``Complete''
means deductively complete in the displayed language; conditional negative
entries expose their consistency assumptions.}
\label{tab:p03-first-order-theories}\\
\toprule
Theory and exact language & Complete? & Decidable? & Literal QE? & Categoricity summary & Exact scope or caveat \\
\midrule\endfirsthead
\toprule
Theory and exact language & Complete? & Decidable? & Literal QE? & Categoricity summary & Exact scope or caveat \\
\midrule\endhead
\(\mathsf{DLO}\) in \(\{<\}\) & Yes & Yes & Yes & \(\aleph_0\)-categorical & Not categorical in any uncountable cardinal. \\
\addlinespace
\(\mathsf{ACF}_p\), fixed \(p=0\) or prime, in \(\{0,1,+,-,\mathbin{\cdot}\}\) & Yes & Yes & Yes & Uncountably categorical & Categorical in every uncountable cardinal; not \(\aleph_0\)-categorical. \\
\addlinespace
\(\mathsf{RCF}\) in \(\{0,1,+,-,\mathbin{\cdot},<\}\) & Yes & Yes & Yes & No infinite categoricity & The order symbol is essential to the literal QE entry. \\
\addlinespace
\(\operatorname{Th}(\N;0,1,+,<)\) (Presburger arithmetic) & Yes & Yes & No & No infinite categoricity & QE holds after adding all congruence, equivalently divisibility, predicates. \\
\addlinespace
\(\operatorname{Th}(\N;S)\) in \(\{S\}\) & Yes & Yes & No & Uncountably categorical & Not \(\aleph_0\)-categorical; naming the unique initial point yields QE. \\
\addlinespace
Atomless Boolean algebras in the Boolean-algebra signature just fixed & Yes & Yes & Yes & \(\aleph_0\)-categorical & Not categorical in uncountable cardinals. \\
\addlinespace
\(\mathsf{PA}\) in \(\{0,S,+,\times\}\) & No, if consistent & No, if consistent & No, if consistent & No infinite categoricity, if consistent & Effective axiomatizability and arithmetic strength activate incompleteness. \\
\addlinespace
True arithmetic \(\operatorname{Th}(\N;0,1,+,\times,<)\) & Yes, by definition & No; not c.e. & No & No infinite categoricity & It has nonstandard elementary equivalents; completeness is not categoricity. \\
\addlinespace
\(\mathsf{ZFC}\) in \(\{\in\}\) & No, if consistent & No, if consistent & No, if consistent & No infinite categoricity, conditional on consistency & The row presupposes that the consistency/model condition is met. \\
\bottomrule
\end{longtable}
\endgroup
\end{landscape}

The DLO, fixed-characteristic ACF, and RCF classifications are standard
consequences of their elimination theorems; the arithmetic negative rows use
the effective and incompleteness boundaries developed elsewhere in the
handbook.  Sources for the table include
  \cite{Marker2002,Steinitz1910,Hodges1993,Presburger1929,Shoenfield1967,Godel1931,
FerranteRackoff1979,Kozen1980}.  No row infers categoricity merely from
completeness.

\begin{landscape}
\begingroup\tiny\setlength{\tabcolsep}{2pt}
\begin{longtable}{@{}p{.65in}p{1.10in}p{1.20in}p{.60in}p{.70in}p{.80in}p{2.10in}@{}}
\caption{Quantifier-elimination catalog.  An em dash in the
definable-choice column records that this entry deliberately makes no claim;
it does not mean failure.}
\label{tab:p03-qe-catalog}\\
\toprule
Theory & Base language and literal QE & Eliminating expansion or witness & Model complete? & Definable Skolem functions? & Complete / decidable? & Calibrated complexity statement \\
\midrule\endfirsthead
\toprule
Theory & Base language and literal QE & Eliminating expansion or witness & Model complete? & Definable Skolem functions? & Complete / decidable? & Calibrated complexity statement \\
\midrule\endhead
\(\mathsf{DLO}\) & \(\{<\}\): yes & Finite lower/upper-bound order types & Yes & No: automorphisms fixing \(x\) prevent a definable choice above \(x\) & Yes / yes & Validity is \(\mathsf{PSPACE}\)-complete \cite{Kozen2006}. \\
\addlinespace
\(\mathsf{ACF}_p\) & Ring language, fixed \(p\): yes & Algebraic-geometric elimination; resultants witness the procedure & Yes & No: conjugate roots cannot be uniformly selected & Yes / yes & Grigor'ev's bound is polynomial in \(L^{n^{2a+1}}\), parameterized by formula length \(L\), variables \(n\), and alternations \(a\); no unqualified polynomial-time claim \cite{Grigoriev1987}. \\
\addlinespace
\(\mathsf{RCF}\) & Ordered-ring language: yes; pure-ring language: no literal claim & Sign conditions and Tarski--Seidenberg projection & Yes & Yes: semialgebraic definable choice & Yes / yes & Some QE outputs require doubly exponential time and space; this is not an exact decision-class classification \cite{DavenportHeintz1988}. \\
\addlinespace
Presburger & \(\{0,1,+,<\}\): no & Add congruence/divisibility predicates; Cooper-style elimination & Yes & --- & Yes / yes & Fischer--Rabin give a doubly exponential lower bound in formula length in their computation model \cite{FischerRabin1974}. \\
\addlinespace
\(\operatorname{Th}(\N;S)\) & \(\{S\}\): no & Add the initial-point constant \(0\); predecessor existence becomes \(x\ne0\) & No in the base language; yes after naming \(0\) & --- & Yes / yes & The one-successor theory in \(\{S\}\) is \(\mathsf{PSPACE}\)-complete \cite{FerranteRackoff1979}. \\
\addlinespace
Atomless Boolean algebras & Boolean-algebra signature: yes & Split finite Boolean combinations into their Boolean atoms & Yes & --- & Yes / yes & Complete under log-space reductions for \(\bigcup_{c<\omega}\mathrm{STA}(*,2^{cn},n)\), including the atomless theory; do not relabel this class \(\mathsf{PSPACE}\) \cite{Kozen1980}. \\
\bottomrule
\end{longtable}
\endgroup
\end{landscape}

The shift embedding \(n\mapsto n+1\) maps \((\N,S)\) onto a substructure
whose first element has no predecessor internally but does have one in the
ambient structure; hence the bare successor theory is not model complete.
Naming \(0\) blocks that embedding, and quantifier elimination in
\(\{0,S\}\) gives model completeness.  The three em dashes in the choice
column keep that independent property open rather than manufacturing it from
QE.

%% END AUDITED QE BODY SLICE 104-212

%% BEGIN AUDITED QE BODY SLICE 326-329
\begin{hbtheorem}{thm:p03-vaught-two-cardinal}{cited=Vaught1965CardinalsFarApart}{Vaught's two-cardinal theorem}
Let \(T\) be a countable first-order theory in a language containing a unary
predicate \(U\). If \(T\) has a model \(M\) with
\(|M|>|U^M|\ge\aleph_0\), then \(T\) has a model \(N\) with
\(|N|=\aleph_1\) and \(|U^N|=\aleph_0\). The countability of the language and
the distinguished unary predicate are part of the statement; the theorem does
not prescribe arbitrary pairs of infinite cardinals.
\end{hbtheorem}

\noindent\emph{Method sketch.}
Expand the theory so that a long indiscernible sequence lies outside the
distinguished predicate while Skolem closure over countably many parameters
keeps the predicate part countable. Compactness supplies finite pieces, and an
Ehrenfeucht--Mostowski construction stretches the sequence to length
\(\aleph_1\). This records the mechanism but not every Skolemization and
indiscernibility lemma; the theorem's proof status remains the cited one.
%% END AUDITED QE BODY SLICE 326-329

%% BEGIN AUDITED QE BODY SLICE 331-451
\section{Ehrenfeucht--Fra\"iss\'e games}
\label{sec:p03-ef-games}

\begin{hbdefinition}{def:p03-quantifier-rank}{Quantifier rank and bounded agreement}
The \emph{quantifier rank} \(\operatorname{qr}(\varphi)\) is defined
recursively: atomic formulas and \(\bottom\) have rank \(0\); negation preserves rank; a
binary Boolean connective takes the maximum rank of its arguments; and
\(\exists x\,\varphi\) and \(\forall x\,\varphi\) have rank
\(\operatorname{qr}(\varphi)+1\).  Two structures \emph{agree through rank
\(k\)} if they satisfy the same first-order sentences of quantifier rank at
most \(k\).  This metalinguistic relation is needed to state the finite-round
game theorem and is weaker than elementary equivalence, which quantifies over
sentences of every finite rank.
\index{quantifier rank}
\end{hbdefinition}

\begin{hbdefinition}{def:p03-ef-game}{Ehrenfeucht--Fra\"iss\'e game}
Fix a finite relational signature \(\tau\), two \(\tau\)-structures
\(\mathcal A,\mathcal B\), and \(k\in\N\).  In each of \(k\) rounds,
Spoiler chooses an element in either structure and Duplicator responds with an
element of the other.  Record every pair as \(a_i\mapsto b_i\), regardless of
which structure Spoiler chose from.  Duplicator wins exactly when the resulting
pairing is a finite partial isomorphism: it respects repetitions/equality and
preserves and reflects every relation symbol on selected tuples.  The finite
relational vocabulary is fixed; function symbols require a separate
relationalization or closure formulation.
\index{Ehrenfeucht--Fraisse game@Ehrenfeucht--Fra\"iss\'e game}
\end{hbdefinition}

\begin{hbtheorem}{thm:p03-ef}{cited=Ehrenfeucht1961}{Ehrenfeucht--Fra\"iss\'e theorem}
For structures \(\mathcal A\) and \(\mathcal B\) in the same fixed finite
relational signature, Duplicator has a winning strategy in the \(k\)-round
Ehrenfeucht--Fra\"iss\'e game if and only if \(\mathcal A\) and
\(\mathcal B\) agree through rank \(k\).
\end{hbtheorem}

The theorem converts syntactic rank into a combinatorial invariant; it does
not turn one fixed finite game into full elementary equivalence.  The game,
rank induction, and its finite-model applications are developed in
\cite{Ehrenfeucht1961,Libkin2004}.

\begin{hbcorollary}{thm:p03-ef-inexpressibility}{cited=Libkin2004}{Game criterion}
Let \(\mathcal C\) be an isomorphism-closed class of finite structures in one
fixed finite relational signature.  If, for every \(k\), there are finite
\(\mathcal A_k\in\mathcal C\) and \(\mathcal B_k\notin\mathcal C\) on which
Duplicator wins the \(k\)-round game, then no first-order sentence defines
\(\mathcal C\) on finite structures.
\end{hbcorollary}

\begin{figure}[htbp]
\centering
\begin{tikzpicture}[
  x=1.35cm,y=1.25cm,
  point/.style={circle,draw,minimum size=6.5mm,inner sep=1pt,font=\small},
  chosen/.style={circle,draw,double,minimum size=7.5mm,inner sep=1pt,font=\small},
  order/.style={-{Latex[length=2mm]}},
  pair/.style={dashed,thick}
]
  \node[point]  (a0) at (0,1.7) {\(a_0\)};
  \node[chosen] (a1) at (1,1.7) {\(a_1\)};
  \node[chosen] (a2) at (2,1.7) {\(a_2\)};
  \draw[order] (a0) -- (a1);
  \draw[order] (a1) -- (a2);
  \node[left=5mm of a0,font=\small] {\(\mathcal A\)};

  \node[point]  (b0) at (-.5,0) {\(b_0\)};
  \node[chosen] (b1) at (.5,0) {\(b_1\)};
  \node[point]  (b2) at (1.5,0) {\(b_2\)};
  \node[chosen] (b3) at (2.5,0) {\(b_3\)};
  \draw[order] (b0) -- (b1);
  \draw[order] (b1) -- (b2);
  \draw[order] (b2) -- (b3);
  \node[left=5mm of b0,font=\small] {\(\mathcal B\)};

  \draw[pair] (a1) -- node[left,font=\footnotesize] {round 1} (b1);
  \draw[pair] (a2) -- node[right,font=\footnotesize] {round 2} (b3);
\end{tikzpicture}
\caption{A two-round EF position on the three- and four-element linear orders.
Spoiler first chose \(a_1\), answered by \(b_1\), and
then chose \(b_3\), answered by \(a_2\).  The selected pairing preserves
equality and order.  This board illustrates one play of a fixed-round strategy;
it is not the uniform family needed for the parity result.}
\label{fig:p03-ef-board}
\end{figure}

For the full two-round strategy in Figure~\ref{fig:p03-ef-board}, Duplicator
answers an endpoint by the corresponding endpoint and an interior first move
by an interior point.  On round two, the response repeats the paired element
when Spoiler repeats a move, and otherwise lies on the same side of the first
selected point.  With only two selected points this preserves their complete
order type.

\begin{hbexample}{def:p03-ef-parity}{Finite order parity is not first-order definable}
For \(k\geq1\), compare the linear orders of sizes \(2^k\) and \(2^k+1\).
With \(r\) rounds remaining, Duplicator maintains that the selected pairing is
order-preserving and every pair of corresponding gaps either has the same size
below \(2^r-1\), or both gaps have size at least \(2^r-1\).  Initially the
invariant holds.  A repeated selection receives its existing mate.  In equal
small gaps Duplicator copies the exact offset.
In two large gaps, if Spoiler leaves fewer than \(2^{r-1}-1\) points on one
side, Duplicator copies that smaller side; otherwise Duplicator chooses a point
leaving at least \(2^{r-1}-1\) points on both sides.  Thus the invariant
survives, and at \(r=0\) the selected map preserves equality and order.  The
\(k=0\) game is vacuous.  Hence every rank has an even/odd pair that agrees
through that rank, so Corollary~\ref{thm:p03-ef-inexpressibility} shows that no
first-order sentence defines even cardinality on finite linear orders.
\end{hbexample}

\begin{hbexample}{def:p03-ef-connectivity}{Finite graph connectivity is not first-order definable}
For each \(k\), let \(C_m\) denote the simple undirected cycle on \(m\)
vertices, and compare \(C_{2m}\) with the disjoint union \(C_m\sqcup C_m\)
for sufficiently large \(m\).  The cited EF argument matches the bounded
neighborhoods of selected vertices: with \(r\) rounds left it preserves all
distances below \(2^r\), while points outside those neighborhoods are answered
outside the corresponding neighborhoods.  For sufficiently large \(m\), the
responses remain available through \(k\) rounds.  One graph is connected and
the other is not, so connectivity is not first-order definable on finite
graphs.  The full invariant and response lemma are cited from
\cite{Libkin2004}; the preceding parity example supplies the fully expanded
strategy pattern used by this entry.
\end{hbexample}
%% END AUDITED QE BODY SLICE 331-451
